\scanpage{355}
% @ head of sheet
\mspage{79}
% @ top, lines 1–3
$X_j\supset X_i$ et $Y_j\supset Y_i$ de cardinal $\leq\unsure{\kappa|\pi}$, telles que $f$ induise $f_j:X_j\unsure{\mathbin{-}|\to}Y_j$ et que l'on ait $f_j\in W$.

% @ upper middle, lines 4–11
La notion d'\unsure{accessibilité} de $W$ me semble aussi bien comprise que celle des limites, je n'y vois plus aucun mystère. D'autre part, il me semble que tous les localiseurs fondamentaux utilisés à ce jour y satisfont, (tout comme à l'axiome $\mathrm{Loc}\,8$). \deleted{Pour} Ici, je vais me borner à établir :

% @ lower middle, theorem heading and statement
\uline{\deleted{Proposition}} \inserted{Théorème} 5.4. Le localiseur \unsure{fond} $\mathbb W_\infty$ (équivalences faibles ``\unsure{habituelles|traditionnelles}) satisfait à l'axiome $\mathrm{Loc}(9)$ (en fait, il satisfait à \uline{tous} les axiomes \unsure{vus|vu} jusqu'\unsure{à} présent, savoir $\mathrm{Loc}(1\mathrm{ter},2,3\mathrm{bis},4,5,6,7,8,9)$).
% @ lower left, sloping note beside Theorem 5.4
% slanted
\margin{Mérite bien le nom de \uline{théorème}, surtout si on y inclut \deleted{\unsure{toute}} la liste des axiomes $\mathrm{Loc}$ mentionnés ci-contre ! Ce théorème n'a rien \deleted{\unsure{de}} d'\unsure{évident}, à partir de la définition \unsure{charabialgique|charabialogique|charabiologique} de $\mathbb W_\infty$ -- c'est un résultat (``bien connu'') qui me paraît profond.}
% @ lower third, small lines below the axiom list
Plus précisément, $\mathbb W_\infty$ est $\unsure{\aleph|\mathbb N}_0$-accessible (où $\aleph_0$ est le premier cardinal infini).
% @ foot, last three lines
Je ne reviens pas \uline{ici} sur la vérification des autres axiomes, et me borne à prouver $\mathrm{Loc}(9)$, l'accessibilité, en utilisant ce qui peut m'être utile. Je me
% archive sheet number: 354
